\documentclass[instructions.tex]{subfiles}
\begin{document}
 
\chapter{Des mauvais livres.}
 
\primo Si l'on faisait contre l'intérêt d'un royaume ce que l'on fait contre la religion, il serait bientôt détruit. Le démon n'a rien inventé de plus efficace pour corrompre les mœurs et la foi, que les mauvais livres. Si la religion n'était pas l'ouvrage de Dieu, elle serait déjà détruite par ce moyen pernicieux.
 
Des mauvais livres, les uns inspirent l'esprit du monde, l'impureté : tels sont certains romans, les histoires d'intrigues galantes, les poësies et les comédies obscènes. Les autres inspirent l'erreur, l'impiété et l'athéisme : tels sont les livres écrits contre l'Eglise romaine, les satires contre la religion, les libelles remplis de calomnies contre les personnes sacrées.
 
Les personnes de condition, étant plus en état de soutenir la religion par leurs exemples, ayant plus d'autorité et plus d'éducation, doivent aussi être plus en garde contre l'ennemi du salut. Il leur tend des pièges plus séduisans pour les pervertir, et y réussit par les mauvais livres.
 
\secundo Quelle horreur n'aurait-on pas de ces livres abominables, si l'on en connaissait les auteurs, la plupart apostats, sans religion, sans pudeur, gens proscrits, dignes de la roue et du feu ! Comment des personnes d'honneur peuvent-elles faire leur occupation de telles lectures?
 
Vous ne voudriez pas meubler vos appartemens des choses qui auraient appartenu à des gens infâmes, et vous en remplissez votre esprit ; vous meublez votre mémoire des productions honteuses de ces auteurs détestables. Vous auriez horreur d'entrer en commerce avec le démon; cependant vous le faites, lorsque vous lisez des livres qui sont les organes de Satan. Dieu nous parle et nous instruit par les bons livres, dit un saint Père, et le démon parle et séduit par les mauvais. Vous auriez honte de faire instruire vos enfans par des gens décriés et sans honneur, et vous vous instruisez vous-même à votre perte, par leurs ouvrages et par leurs livres.
 
Si un écrit séditieux contre le roi et le bon ordre de l'état tombait entre vos mains, loin de le lire et de le communiquer, vous le condamneriez au feu ; et vous n'avez point d'horreur de communiquer aux autres et de lire des libelles infâmes, écrits contre ce que vous avez de plus sacré et de plus cher au monde, qui est la religion et la pureté des mœurs !
 
\tertio Nous lisons ces livres, dit-on, pour nous former, et pour apprendre la langue dans sa pureté. Mais n'y a-t-il pas de bons livres pour vous former l'esprit? Par la lecture des mauvais, dit saint Augustin, on n'apprend pas à devenir éloquent, mais à devenir vicieux; on y apprend à connaître le mal sans horreur, à en parler sans pudeur, à le commettre sans retenue. On veut, par ces lectures, se former l'esprit, et l'on s'y pervertit. On y perd la droiture du jugement, on y apprend à être pointilleux, impudent, incrédule et athée. De là vient que certaines gens, qui raisonnent en hommes sensés sur les affaires du monde, ne raisonnent qu'en pitoyables sophistes sur la religion. Les livres qui les ont séduits n'étant remplis que de fausses suppositions, de faux principes, de faux raisonnemens, et d'un faux brillant, dès qu'on les a goûtés, c'est un mal presque incurable. On a en horreur tous les bons livres, on n'est même presque plus capable de raisonnement en matière de religion et de bonnes mœurs.
 
Étrange bizarrerie de ces esprits égarés ! Trouvent-ils dans un bon livre quelques faits merveilleux et édifians, ils n'en croient rien. Trouvent-ils dans un mauvais livre des impertinences, des faussetés, des faits supposés contre l'Eglise, ils les croient. Rencontrent-ils quelques réflexions solides sur l'autre vie, sur les maximes de l'Evangile, dans un livre de piété, ils s'en dégoûtent. Trouvent-ils quelques fades plaisanteries contre la pudeur, quelques traits ridicules contre la religion, dans un libelle, ils le goûtent et le dévorent. C'est donc ainsi que Dieu abandonne ces esprits orgueilleux à leur sens réprouvé. O aveuglement qui leur fait prendre la vérité pour le mensonge, et le mensonge pour la vérité !
 
Si un homme avait le sang corrompu, le verrait-on, au mépris des médecins, faire usage du poison pour se rétablir ? Pourquoi, vous qui avez le cœur gâté par tant de passions, l'empoisonnez-vous encore par ces lectures envenimées, au mépris des livres de piété, et des pasteurs qui sont destinés à vous instruire ?
 
Plus un mauvais livre vous paraît agréable, rempli de traits délicats et éblouissans, plus il est pernicieux; plus le poison en est doux, et plus il est dangereux. Un mauvais livre est le plus cruel ennemi que vous ayez dans votre maison : condamnez-le au feu, il ne mérite pas une autre destinée. Si vous le gardez, c'est une vipère qui tôt ou tard vous fera des blessures mortelles.
 
Malheur à ceux qui composent de tels ouvrages! Malheur à ceux qui les impriment, qui les autorisent, qui les vendent, qui les distribuent, qui les prêtent! Avec quelle force les magistrats emploieraient-ils leur autorité pour supprimer un écrit contre le prince, et en punir les auteurs! L'intérêt de Dieu est le seul pour lequel on manque de zèle. On voit une foule de livres contre l'Eglise, contre la religion, contre la pureté des mœurs; les magistrats s'endorment; les pères, les maîtres les permettent, les pasteurs se récrient, et on les méprise. On devrait verser des larmes de sang sur de tels abus.
 
\section{Résolutions}
 
\primo Je regarderai toute ma vie les mauvais livres comme un poison dont il ne faut pas infecter son esprit, si on veut sauver son âme. 

\secundo Je veillerai soigneusement sur les personnes qui dépendent de moi, afin de les empêcher de lire rien qui puisse corrompre leur cœur, gâter leur esprit, et leur inspirer de l'éloignement pour la religion, du mépris pour ses ministres, de l'insoumission envers l'autorité de l'Eglise. 
\tertio En conséquence, si je découvre jamais dans ma famille des romans, des comédies, des livres impies ou hérétiques, je me hâterai de les détruire. 

\prayer{O mon Dieu ! ne permettez pas que ni moi, ni ceux de qui je dois répondre devant vous, nous nous exposions à nous perdre pour l'éternité par de mauvaises lectures.}
 
\end{document}
